\documentclass[10pt]{article}

\usepackage[margin=0.8in]{geometry}
\usepackage{amsmath,amssymb,amsthm,mathtools}
\usepackage{booktabs}
\usepackage{microtype}
\usepackage[hidelinks]{hyperref}
\usepackage{xurl}
\emergencystretch=2em

\newtheorem{theorem}{Theorem}
\newtheorem{proposition}[theorem]{Proposition}
\newtheorem{lemma}[theorem]{Lemma}
\newtheorem{corollary}[theorem]{Corollary}
\theoremstyle{remark}
\newtheorem{remark}[theorem]{Remark}

\newcommand{\F}{\mathbb F}
\newcommand{\rank}{\operatorname{rank}}
\newcommand{\supp}{\operatorname{supp}}
\newcommand{\GL}{\operatorname{GL}}

\hypersetup{
  pdftitle={The Exact F2-Rank of 2x3 by 3x4 Matrix Multiplication},
  pdfauthor={Ian D'Ambrosio}
}

\title{The Exact $\F_2$-Rank of $2\times3$ by $3\times4$\
Matrix Multiplication}
\author{Ian D'Ambrosio\\
\small Nth Research Collective\\
\small \texttt{ian@nthresearch.org}}
\date{11 August 2026}

\begin{document}
\maketitle

\begin{abstract}
Hopcroft and Kerr gave a 20-multiplication algorithm for multiplying a
$2\times3$ matrix by a $3\times4$ matrix, and conjectured that 20
multiplications are necessary.  Over $\F_2$, Wang recently proved the lower
bound 19 with a machine-checkable substitution certificate, leaving the exact
rank in $\{19,20\}$.  We prove that the rank is 20.  Wang's checked
first-factor capacities imply that every hypothetical 19-term decomposition
has pairwise distinct first factors and at most one first factor of rank two.  A
checked CNF/DRAT/LRAT certificate and exhaustive enumeration reduce all such
supports to 252 profiles in four symmetry orbits.  Row restrictions then
exclude the four orbits.  Tight 12-term restrictions force output factors into
a four-dimensional row space; 13-term restrictions have a unique relation
whose quotient factors are all equal.  These facts immediately exclude the
all-rank-one profiles of types $(7,7,5)$ and $(7,6,6)$.  In the remaining
profile, with type $(6,6,6)$ and one rank-two factor, the three relations would
need at least 18 support incidences but can contain at most 9.  An explicit
20-term decomposition is checked against all 576 tensor coefficients.  The
profile proof, restriction certificate, upper witness, and replay commands are
provided as exact artifacts.
\end{abstract}

\noindent\textbf{Keywords.}
bilinear complexity; tensor rank; matrix multiplication; finite fields;
computer-assisted proof

\section{Introduction}

Let $R_{\F}(\langle m,n,p\rangle)$ denote the tensor rank over $\F$ of
multiplying an $m\times n$ matrix by an $n\times p$ matrix.  Equivalently, it
is the least number of scalar multiplications in a bilinear algorithm for this
map.  Exact ranks for small rectangular formats are basic finite problems in
algebraic complexity, but even tiny cases can depend on the ground field and
resist both general lower bounds and constructive search
\cite{Blaser2003,Smirnov2013}.

For $\langle2,3,4\rangle$, Hopcroft and Kerr gave an algorithm using 20
multiplications and conjectured that 20 are necessary
\cite{HopcroftKerr1971}.  Smirnov's later table distinguishes an exact upper
bound 20 from an approximate bound 19 \cite{Smirnov2013}.  AlphaTensor also
found an exact rank-20 decomposition over $\F_2$ \cite{FawziEtAl2022}.
Wang's automated finite-field framework recently supplied the first checked
lower bound
\[
 R_{\F_2}(\langle2,3,4\rangle)\geq19,
\]
while retaining 20 as the upper bound \cite{Wang2026}.  Thus only one exact
case remained.

\begin{theorem}\label{thm:main}
The bilinear complexity over $\F_2$ of multiplying a $2\times3$ matrix by a
$3\times4$ matrix is
\[
 R_{\F_2}(\langle2,3,4\rangle)=20.
\]
\end{theorem}

The lower-bound proof has a certified finite front end and a short structural
back end.  Wang's substitution bounds reduce any hypothetical rank-19
decomposition to four orbits of possible first-factor supports.  We then use
the three nonzero row restrictions of $\F_2^2$.  A restriction has flattening
rank 12.  With 12 surviving terms it is tight and rigid; with 13 terms it has
exactly one relation.  The resulting contradictions depend only on the orbit's
row multiplicities, so no search over the remaining two factors is needed.

\paragraph{Contribution and scope.}
The result closes this exact $\F_2$ format and supplies a new restriction-based
exclusion of rank 19.  The argument uses the special three-direction geometry
of $\F_2^2$ and does not claim a general rank formula for
$\langle2,3,p\rangle$ or for other fields.  Searches over flip graphs and
fixed-profile SAT instances were useful diagnostics but play no role in the
proof.

\section{Tensor conventions and row restrictions}

Put
\[
 A=\F_2^2,\qquad B=\F_2^3,\qquad C=\F_2^4.
\]
Using the standard bases to identify coordinate spaces with their duals, the
matrix-multiplication tensor is
\begin{equation}\label{eq:tensor}
 M_{2,3,4}
 =\sum_{i=1}^{2}\sum_{j=1}^{3}\sum_{k=1}^{4}
   a_{ij}\otimes b_{jk}\otimes c_{ik}
 \in (A\otimes B)\otimes(B\otimes C)\otimes(A\otimes C).
\end{equation}
Suppose, toward a contradiction, that
\begin{equation}\label{eq:rank19}
 M_{2,3,4}=\sum_{s=1}^{19}U_s\otimes V_s\otimes W_s,
\end{equation}
where every factor is nonzero.  Regard $U_s$ as a $2\times3$ binary matrix.

For $x\in A\setminus\{0\}$, restrict the left input to matrices $xy^{\mathsf
T}$ with $y\in B$.  Define
\[
 H_x=x\otimes C\subseteq A\otimes C,
 \qquad R_x=B\otimes H_x.
\]
The restricted target is
\begin{equation}\label{eq:restricted}
 T_x=\sum_{j=1}^{3}\sum_{k=1}^{4}
 e_j\otimes E_{jk}\otimes(x\otimes f_k).
\end{equation}
Its flattening with $B\otimes C$ on one side has rank 12 and image exactly
$R_x$.  Moreover,
\begin{equation}\label{eq:intersection}
 H_x\cap H_y=\{0\}\qquad(x\neq y,
 \ x,y\in A\setminus\{0\}).
\end{equation}

If $U_s=a_sb_s^{\mathsf T}$ has rank one, then term $s$ survives restriction
$x$ exactly when $a_s\mathbin{\cdot}x=1$, and its second flattening factor is
\begin{equation}\label{eq:q-rank1}
 q_{s,x}=b_s\otimes W_s.
\end{equation}
If $U_*$ has rank two, it survives all three nonzero restrictions and
\begin{equation}\label{eq:q-rank2}
 q_{*,x}=(x^{\mathsf T}U_*)\otimes W_*.
\end{equation}
The first factor in \eqref{eq:q-rank2} is nonzero because $U_*$ has full row
rank.

\section{Two flattening lemmas}

\begin{lemma}[Tight restriction]\label{lem:tight}
If $T_x$ is expressed by exactly 12 active terms from
\eqref{eq:rank19}, then the active $V_s$ form a basis of $B\otimes C$ and the
active $q_{s,x}$ form a basis of $R_x$.  In particular, every active rank-one
term satisfies $W_s\in H_x$.
\end{lemma}

\begin{proof}
A rank-12 matrix written as a sum of 12 dyads forces both factor lists to have
rank 12.  Hence the $q_{s,x}$ span, and therefore belong to, the flattening
image $R_x$.  If $b_s\otimes W_s\in B\otimes H_x$ with $b_s\neq0$, apply a
functional taking $b_s$ to 1 to obtain $W_s\in H_x$.
\end{proof}

\begin{lemma}[One-excess restriction]\label{lem:excess}
Suppose that $T_x$ is expressed by 13 active terms, indexed by $I$.  Let
\[
 L_x:\F_2^I\longrightarrow B\otimes C,
 \qquad e_s\longmapsto V_s.
\]
Then $L_x$ is surjective and
$\ker L_x=\langle\lambda^x\rangle$ is one-dimensional.  If
\[
 \pi_x:B\otimes(A\otimes C)\longrightarrow
 \bigl(B\otimes(A\otimes C)\bigr)/R_x
\]
is the quotient map, there is a quotient vector $z_x$ such that
\begin{equation}\label{eq:quotient-relation}
 \pi_x(q_{s,x})=\lambda_s^x z_x\qquad(s\in I).
\end{equation}
\end{lemma}

\begin{proof}
The active $V_s$ span 12 dimensions because the restricted flattening has rank
12; their ambient space also has dimension 12.  Thus $L_x$ is surjective and
has a one-dimensional kernel.  Projecting the target outside $R_x$ gives
\[
 \sum_{s\in I}V_s\otimes\pi_x(q_{s,x})=0.
\]
Equivalently, the tensor
$\sum_{s\in I}e_s\otimes\pi_x(q_{s,x})$ belongs to
$\ker L_x\otimes\bigl((B\otimes(A\otimes C))/R_x\bigr)$.  Since the first
factor is spanned by $\lambda^x$, equation
\eqref{eq:quotient-relation} follows.
\end{proof}

We will also use that over $\F_2$, equality of nonzero pure tensors
$b\otimes h=b'\otimes h'$ forces $b=b'$ and $h=h'$.  Over a general field the
factors are proportional; $\F_2$ has only one nonzero scalar.

\section{The certified first-factor profiles}

Wang's substitution certificate supplies a checked residual lower bound
$L(S)$ for each relevant subspace $S$ of the first tensor factor.  The bound
applies after restricting the left input to the annihilator $S^\perp$, which
kills exactly the terms whose first factor belongs to $S$.  If $m(S)$ terms of
a rank-19 decomposition have $U_s\in S$, substitution gives the capacity
inequality
\begin{equation}\label{eq:capacity}
 m(S)+L(S)\leq19.
\end{equation}
For every one-dimensional $S$, the checked residual bound is 18.  Since a
one-dimensional binary subspace has one nonzero vector, every nonzero first
factor therefore occurs at most once.

The profile certificate expands Wang's orbit data to all first-factor
subspaces and retains the 1,897 nontrivial inequalities
\eqref{eq:capacity}.  There are 21 nonzero $2\times3$ matrices of rank one and
42 of rank two.  A 62,267-variable CNF asserts all capacities, cardinality 19,
and at least two selected rank-two matrices.  Its 124,027 clauses are
unsatisfiable; both its DRAT and LRAT proofs check.  Exhaustive enumeration of
the remaining zero- or one-rank-two candidates checks 56,070 supports and
retains exactly 252.  Quotienting by
$\GL(2,2)\times\GL(3,2)$, of order 1,008, gives the following complete list.

\begin{proposition}[Checked profile classification]\label{prop:profiles}
The first factors in any rank-19 decomposition \eqref{eq:rank19} are pairwise
distinct, at most one has matrix rank two, and their support belongs to one of
the four orbits in Table~\ref{tab:profiles}.
\end{proposition}

\begin{table}[ht]
\centering
\caption{The four certified first-factor support orbits.  The type records the
multiplicities of the three nonzero row directions among rank-one factors.}
\label{tab:profiles}
\begin{tabular}{@{}cccl@{}}
\toprule
orbit & size & type & rank-two factors \\
\midrule
0 & 63  & $(7,7,5)$ & none \\
1 & 21  & $(7,6,6)$ & none \\
2 & 126 & $(7,6,6)$ & none \\
3 & 42  & $(6,6,6)$ & one \\
\bottomrule
\end{tabular}
\end{table}

The CNF, proof traces, enumeration, group action, and orbit representatives are
reconstructed by the verifier described in Section~\ref{sec:verification}.
Proposition~\ref{prop:profiles} is the only computer-assisted classification
used below.

\section{Excluding the four profiles}

\subsection{\texorpdfstring{Type $(7,7,5)$}{Type (7,7,5)}}

Two row directions occur seven times.  Restricting against the annihilator of
either direction removes those seven terms and leaves exactly 12.  By
Lemma~\ref{lem:tight}, every active output factor lies in the corresponding
$H_x$.  The two active sets overlap in the five terms from the deficient row
direction.  Their nonzero output factors would lie in
$H_x\cap H_y=\{0\}$, a contradiction.  Orbit 0 is impossible.

\subsection{\texorpdfstring{Type $(7,6,6)$}{Type (7,6,6)}}

Let $a_0$ be the multiplicity-seven direction and let $a_1,a_2$ be the other
two.  Restriction against the annihilator $x_0$ of $a_0$ leaves the 12 terms in
the $a_1$ and $a_2$ groups.  Lemma~\ref{lem:tight} puts all their output
factors in $H_{x_0}$.

Now restrict against the annihilator $x_1$ of $a_1$.  The seven $a_0$ terms
and six $a_2$ terms survive, so Lemma~\ref{lem:excess} applies.  For each of
the six $a_2$ terms, $W_s\in H_{x_0}$, and projection modulo $H_{x_1}$ is
injective on $H_{x_0}$ by \eqref{eq:intersection}.  Its quotient pure tensor
is therefore nonzero.  Equation~\eqref{eq:quotient-relation} forces all six
relevant relation coefficients to be 1 and all six quotient pure tensors to
equal the same nonzero $z_{x_1}$.  The $b_s$ within one row-direction group
are distinct because the $U_s$ are distinct.  Uniqueness of nonzero pure-tensor
factorization gives a contradiction.  Orbits 1 and 2 are impossible.

\subsection{\texorpdfstring{Type $(6,6,6)$ with one rank-two factor}{Type (6,6,6) with one rank-two factor}}

Every nonzero row restriction has 12 active rank-one terms plus the rank-two
term.  Lemma~\ref{lem:excess} applies to all three restrictions.

First, $z_x\neq0$ for every nonzero $x$.  Otherwise
\eqref{eq:quotient-relation} places every output factor active under $x$ in
$H_x$.  Choose $y\neq x$.  The two active sets overlap in one complete
six-term rank-one row group and the rank-two term.  The seven overlap output
factors are nonzero and lie in $H_x$, so none lies in $H_y$.  Their quotient
tensors under restriction $y$ are nonzero.  Equation
\eqref{eq:quotient-relation} would make the six rank-one quotient tensors equal,
contradicting their distinct $b_s$ factors.

Since $z_x\neq0$, every term in $\supp(\lambda^x)$ has the same nonzero
quotient pure tensor.  At most one term can come from either active rank-one
row group, and at most one is the rank-two term.  Hence
\begin{equation}\label{eq:support-capacity}
 |\supp(\lambda^x)|\leq3
\end{equation}
for each of the three restrictions.

Each of the 18 rank-one terms is active under exactly two restrictions.  If
its relation coefficient vanished in both, equation
\eqref{eq:quotient-relation} would put its nonzero output factor in two distinct
spaces $H_x$ and $H_y$, contradicting \eqref{eq:intersection}.  Thus every one
of the 18 terms occurs in at least one relation support.  The three supports
therefore require at least 18 incidences.  Equation
\eqref{eq:support-capacity} permits at most
\[
 3\text{ restrictions}\times3\text{ terms}=9,
\]
a contradiction.  Orbit 3 is impossible.

\begin{proof}[Proof of Theorem~\ref{thm:main}]
Proposition~\ref{prop:profiles} classifies every first-factor support of a
hypothetical rank-19 decomposition.  The three arguments above exclude all
four orbits, so rank 19 is impossible.  Wang's checked bound excludes ranks at
most 18, while the 20-term witness in Appendix~\ref{app:witness} proves rank at
most 20.  Therefore the rank is exactly 20.
\end{proof}

\section{Exact verification}\label{sec:verification}

All computations use binary or integer arithmetic.  The final certificate
links four independently replayable components:
\begin{enumerate}
\item Wang's source commit and the n324 certificate/backtracking archive; a
      clean Linux build of the pinned verifier reports unconstrained lower
      bound 19 and \texttt{OK};
\item the profile CNF and its checked DRAT and LRAT traces, followed by direct
      enumeration of 56,070 candidates and the four-orbit quotient;
\item the row-restriction packet, regenerated from the 252 profiles and checked
      to exclude every one; and
\item the 20-term witness, expanded into all $6\cdot12\cdot8=576$ tensor
      coefficients and compared with \eqref{eq:tensor}.
\end{enumerate}

From the repository root, the exact-rank replay is
\begin{verbatim}
CERT=output/f2-234-tensor-rank/exact-rank20-certificate-v2.json
python3 -m experiments.f2_234_tight_row_restriction \
  verify-exact-rank20 --artifact "$CERT" --timeout 300
\end{verbatim}
It returns \texttt{EXACT\_RANK\_VERIFIED}, rank 20, 252 excluded rank-19
profiles, and 20 upper-bound terms.  The upper-witness controls and focused
tests run with
\begin{verbatim}
python3 -m evaluators.f2_234_tensor_rank_eval verify-controls
python3 -m pytest tests/test_f2_234_tight_row_restriction.py \
  tests/test_f2_234_profile_orbits.py \
  tests/test_f2_234_tensor_rank_eval.py -q
\end{verbatim}
The controls accept the rank-20 witness and Strassen's rank-7 tensor, and reject
a one-bit corruption, a zero factor, and a false declared rank.  The focused
suite passes 35 tests.

The final certificate has SHA-256 digest
\begin{center}
\footnotesize\ttfamily
89dbb5b9de7b13464a7bb7c78d6081a884c9ae933be923a7c787b84fcb5bd29b.
\end{center}
It records the child digests for the profile packet, reviewed exclusion, Wang
log and inputs, upper witness, and profile CNF/DRAT/LRAT proof pair.  A fresh
independent replay re-derived the restriction proof, reproduced the profile
coverage and upper witness, and repeated the primary prior-art search before
publication review.

\paragraph{Data and code availability.}
The exact certificate, proof traces, standalone verifier modules, tests, and
paper source form the reproducibility package accompanying this manuscript.  A
permanent repository identifier will be added upon archival deposit; no result
in this paper depends on that future deposit.

\paragraph{Research-system disclosure.}
Beyond, the research system operated by Nth Research Collective, materially
assisted with literature retrieval, hypothesis generation, exact computation,
proof critique, orchestration, adversarial review, and drafting.  Ian
D'Ambrosio selected the target, directed the investigation, reconciled the
evidence, verified the final claims, and accepts full responsibility for the
manuscript.  Beyond is not an author.

\appendix
\section{An explicit 20-term upper witness}\label{app:witness}

Index $U$ by the six left-input coordinates in row-major order, $V$ by the 12
right-input coordinates in row-major order, and $W$ by the eight output-dual
coordinates in row-major order.  For a binary vector $r=(r_0,\ldots,r_{d-1})$,
write its hexadecimal mask for the integer $\sum_{j=0}^{d-1}r_j2^j$.
Table~\ref{tab:witness} lists the 20 rank-one terms
$u_s\otimes v_s\otimes w_s$.  Expanding the masks gives 24 unit coefficients
at $((i,j),(j,k),(i,k))$ and zero at the other 552 coordinates.

\begin{table}[ht]
\centering
\caption{Hexadecimal factor masks for the exact rank-20 decomposition.}
\label{tab:witness}
\small
\begin{tabular}{@{}rccc@{\qquad}rccc@{}}
\toprule
$s$ & $u_s$ & $v_s$ & $w_s$ & $s$ & $u_s$ & $v_s$ & $w_s$\\
\midrule
1  & \texttt{20} & \texttt{100} & \texttt{90} &
11 & \texttt{20} & \texttt{220} & \texttt{60}\\
2  & \texttt{3e} & \texttt{02a} & \texttt{82} &
12 & \texttt{10} & \texttt{010} & \texttt{11}\\
3  & \texttt{01} & \texttt{008} & \texttt{0a} &
13 & \texttt{12} & \texttt{9d5} & \texttt{01}\\
4  & \texttt{38} & \texttt{020} & \texttt{22} &
14 & \texttt{08} & \texttt{022} & \texttt{f0}\\
5  & \texttt{04} & \texttt{880} & \texttt{09} &
15 & \texttt{16} & \texttt{980} & \texttt{81}\\
6  & \texttt{03} & \texttt{040} & \texttt{44} &
16 & \texttt{13} & \texttt{045} & \texttt{41}\\
7  & \texttt{01} & \texttt{044} & \texttt{05} &
17 & \texttt{3f} & \texttt{62a} & \texttt{02}\\
8  & \texttt{08} & \texttt{023} & \texttt{50} &
18 & \texttt{36} & \texttt{92a} & \texttt{80}\\
9  & \texttt{04} & \texttt{400} & \texttt{06} &
19 & \texttt{06} & \texttt{080} & \texttt{88}\\
10 & \texttt{3b} & \texttt{620} & \texttt{42} &
20 & \texttt{1b} & \texttt{625} & \texttt{40}\\
\bottomrule
\end{tabular}
\end{table}

\begingroup
\small
\begin{thebibliography}{Wang2026}

\bibitem{Blaser2003}
M.~Bl\"aser,
\newblock On the complexity of the multiplication of matrices of small formats,
\newblock \emph{Journal of Complexity} 19 (2003), 43--60.
\newblock \href{https://doi.org/10.1016/S0885-064X(02)00007-9}{doi:10.1016/S0885-064X(02)00007-9}.

\bibitem{FawziEtAl2022}
A.~Fawzi et al.,
\newblock Discovering faster matrix multiplication algorithms with
reinforcement learning,
\newblock \emph{Nature} 610 (2022), 47--53.
\newblock \href{https://doi.org/10.1038/s41586-022-05172-4}{doi:10.1038/s41586-022-05172-4}.

\bibitem{HopcroftKerr1971}
J.~E. Hopcroft and L.~R. Kerr,
\newblock On minimizing the number of multiplications necessary for matrix
multiplication,
\newblock \emph{SIAM Journal on Applied Mathematics} 20 (1971), 30--36.
\newblock \href{https://doi.org/10.1137/0120004}{doi:10.1137/0120004}.

\bibitem{Smirnov2013}
A.~V. Smirnov,
\newblock The bilinear complexity and practical algorithms for matrix
multiplication,
\newblock \emph{Computational Mathematics and Mathematical Physics} 53 (2013),
1781--1795.
\newblock \href{https://doi.org/10.1134/S0965542513120129}{doi:10.1134/S0965542513120129}.

\bibitem{Wang2026}
C.~Wang,
\newblock Automated lower bounds for bilinear complexity over finite fields,
\newblock arXiv:2603.07280v10, 2026.
\newblock \href{https://arxiv.org/abs/2603.07280v10}{arXiv:2603.07280v10}.

\end{thebibliography}
\endgroup

\end{document}
